% ---
% title: A scratch card
% tags: expectation, variance
% ---

\begin{problem}
	A scratch card costs 2 euros. It pays 100 euros with probability $0.001$,
	10 euros with probability $0.02$, 2 euros with probability $0.10$, and
	nothing otherwise. Let $Y$ be the prize (in euros) and let $G=Y-2$ be the
	net gain of a player who buys one card. (a) Find $\E[Y]$ and $\E[G]$.
	(b) Find $\SD(Y)$ and $\SD(G)$.
\end{problem}

\begin{solution}
	\idea Write down the pmf of $Y$ and use $\E[Y]=\sum y\,p(y)$ and
	$\Var(Y)=\E[Y^2]-(\E[Y])^2$. Since $G=Y-2$, the expectation shifts by $-2$
	and the variance does not change.
	\step{The distribution of $Y$} The prize takes four values: $100$, $10$,
	$2$ and $0$. The first three probabilities are given. The card pays exactly
	one of these amounts, so the four probabilities add up to $1$. Hence
	\[ P(Y=0)=1-0.001-0.02-0.10=0.879. \]
	The probability mass function is
	\begin{center}
		\begin{tabular}{c|cccc}
			$y$      & $0$     & $2$    & $10$   & $100$   \\ \hline
			$p(y)$   & $0.879$ & $0.10$ & $0.02$ & $0.001$
		\end{tabular}
	\end{center}
	\step{(a)} Multiply each value by its probability and add:
	\[ \E[Y]=0\cdot0.879+2\cdot0.10+10\cdot0.02+100\cdot0.001
		=0+0.2+0.2+0.1=0.5 \text{ euros}. \]
	The value $0$ contributes nothing, so we may leave it out. The net gain is
	$G=Y-2$: the player always pays 2 euros, whatever the prize. By linearity,
	$\E[aY+b]=a\E[Y]+b$ with $a=1$ and $b=-2$:
	\[ \E[G]=\E[Y]-2=0.5-2=-1.5 \text{ euros}. \]
	On average a player loses 1.5 euros per card.
	\step{(b)} First compute $\E[Y^2]$. Square each value and weight it with the
	same probabilities (the probabilities do not get squared):
	\[ \E[Y^2]=2^2\cdot0.10+10^2\cdot0.02+100^2\cdot0.001
		=0.4+2+10=12.4. \]
	Then
	\[ \Var(Y)=\E[Y^2]-(\E[Y])^2=12.4-0.5^2=12.4-0.25=12.15, \]
	and
	\[ \SD(Y)=\sqrt{12.15}\approx3.486 \text{ euros}. \]
	For $G$ use $\Var(aY+b)=a^2\Var(Y)$ with $a=1$, $b=-2$. Subtracting the
	fixed price moves every possible value by the same amount, so the spread
	does not change:
	\[ \Var(G)=\Var(Y-2)=\Var(Y)=12.15, \qquad \SD(G)=\SD(Y)\approx3.486
		\text{ euros}. \]
	\begin{answer}
		(a) $\E[Y]=0.5$ euros, $\E[G]=-1.5$ euros \quad
		(b) $\SD(Y)=\SD(G)=\sqrt{12.15}\approx3.486$ euros
	\end{answer}
	\begin{remark}
		The expected loss of 1.5 euros is $75\%$ of the 2-euro price: on
		average only a quarter of the money paid for the cards comes back as
		prizes. The standard deviation (about 3.5 euros) is large compared with
		the expected prize of 0.5 euros. This comes from the rare 100-euro
		prize: it contributes only $0.1$ to $\E[Y]$ but $10$ of the $12.4$ in
		$\E[Y^2]$.

		\textbf{Common mistake:} computing $\Var(Y)$ as $\E[Y^2]$ and forgetting
		to subtract $(\E[Y])^2$, or writing $\Var(G)=\Var(Y)-2$. A constant
		shift changes the expectation but never the variance.
	\end{remark}
\end{solution}
